% this file is called up by main.tex
% content in this file will be fed into the main document

% ---------------------------------------------------------------------------

密码学。看起来似乎只有数学家和计算机科学家才能触及这门晦涩、深奥、强大而优雅的学科。事实上，许多密码学的基础概念足够简单，任何人都可以学习。
\par
众所周知，密码学被用于保护通信安全，无论是加密信件还是私密的数字交互。密码学的另一个应用是所谓的加密货币。这些数字货币使用密码学来分配和转移资金所有权。为了确保任何一笔钱都无法被复制或随意创造，加密货币通常依赖于"区块链"——一种公共的、分布式的账本，记录着可以被第三方验证的货币交易 \cite{Nakamoto_bitcoin}。
\par
乍看之下，交易似乎需要以明文形式发送和存储才能被公开验证。事实上，使用密码学工具可以隐藏交易的参与者以及涉及的金额，同时仍然允许观察者验证和确认交易 \cite{cryptoNoteWhitePaper}。加密货币 Monero 就是这一能力的典范。
\par
我们在此努力教授任何懂得基本代数和简单计算机科学概念（如数字的"二进制表示"）的人，不仅要深入全面地理解 Monero 的工作原理，还要领略密码学的实用与优美。
\par
对于有经验的读者：Monero 是一种标准的一维有向无环图（DAG）加密货币区块链 \cite{Nakamoto_bitcoin}，其交易基于椭圆曲线密码学，使用 Ed25519 曲线 \cite{Bernstein2008}，交易输入采用 Schnorr 风格的多层可链接自发匿名群签名（MLSAG）\cite{MRL-0005-ringct} 签名，输出金额（通过 ECDH \cite{Diffie-Hellman} 传达给接收方）使用 Pedersen 承诺 \cite{maxwell-ct} 隐藏，并通过 Bulletproofs \cite{Bulletproofs_paper} 证明其在合法范围内。本报告的前半部分大部分篇幅用于解释这些概念。
